%% Manuscript skeleton on Elsevier's official elsarticle class (v3.5, 2026/01/09).
%% Target: Pattern Recognition or Neurocomputing (docs/DECISIONS.md D-010).
%% Every number comes from generated/results_macros.tex (scripts/make_tables.py);
%% no result is typed by hand. Writing order: Related Work (Phase 0) ->
%% Method/Experiments (Phase 1) -> Results (Phase 2, per supported claim) ->
%% Introduction and Abstract last.
\documentclass[preprint,12pt,authoryear]{elsarticle}

\usepackage[T1]{fontenc}
\usepackage{lmodern}  % Type 1 fonts: vector, embeddable
\usepackage{amsmath,amssymb}
\usepackage{graphicx}
\usepackage{booktabs}
\usepackage{xcolor}
\usepackage[hidelinks]{hyperref}

% AUTO-GENERATED by scripts/make_tables.py from manifests. DO NOT EDIT.
% sources: none yet
\providecommand{\textcolor}[2]{#2}
\newcommand{\DatasetErasedByPolicy}{\textcolor{red}{[pending]}}
\newcommand{\DatasetNumImages}{\textcolor{red}{[pending]}}
\newcommand{\DatasetNumTestDefect}{\textcolor{red}{[pending]}}
\newcommand{\DatasetNumTestGood}{\textcolor{red}{[pending]}}
\newcommand{\DatasetNumTrain}{\textcolor{red}{[pending]}}
\newcommand{\MetricMaxDiffAUPRO}{\textcolor{red}{[pending]}}
\newcommand{\MetricMaxDiffAUROC}{\textcolor{red}{[pending]}}
\newcommand{\SanityConstantAUPRO}{\textcolor{red}{[pending]}}
\newcommand{\SanityConstantImgAUROC}{\textcolor{red}{[pending]}}
\newcommand{\SanityConstantPixAUROC}{\textcolor{red}{[pending]}}
\newcommand{\SanityOracleAUPRO}{\textcolor{red}{[pending]}}
\newcommand{\SanityOracleImgAUROC}{\textcolor{red}{[pending]}}
\newcommand{\SanityOraclePixAUROC}{\textcolor{red}{[pending]}}
\newcommand{\SanityRandomAUPRO}{\textcolor{red}{[pending]}}
\newcommand{\SanityRandomImgAUROC}{\textcolor{red}{[pending]}}
\newcommand{\SanityRandomPixAUROC}{\textcolor{red}{[pending]}}

\newcommand{\todo}[1]{\textcolor{red}{[TODO: #1]}}

\journal{Pattern Recognition}

\begin{document}

\begin{frontmatter}

\title{Does memory compression cost small defects? A controlled study of
normal-feature memory banks in industrial anomaly detection\tnoteref{t1}}
\tnotetext[t1]{Working title. The claim is restricted to what
\texttt{docs/CLAIMS.md} supports at submission.}

\author{\todo{author name}}
\address{\todo{affiliation}}

\begin{abstract}
\todo{Written last, from supported claims only.}
\end{abstract}

\begin{keyword}
industrial anomaly detection \sep memory bank \sep coreset \sep
defect size \sep evaluation protocol \sep MVTec AD
\end{keyword}

\end{frontmatter}

\section{Introduction}\label{sec:intro}
\todo{Written last. Contribution list = the supported rows of CLAIMS.md.}

\section{Related work}\label{sec:related}

\paragraph{Memory-bank anomaly detection and its compression}
Training-free memory-bank detectors score a test patch by its distance to a
bank of normal training-patch features. PatchCore~\citep{roth2022patchcore}
made this the standard industrial baseline. It aggregates locally pooled
mid-level features of an ImageNet-pretrained backbone and reduces the bank
with an approximate greedy $k$-centre coreset. The authors report that the
coreset leaves mean image- and pixel-level performance on MVTec
AD~\citep{bergmann2019mvtec,bergmann2021mvtec} almost unchanged down to a
1\,\% bank. Later work treats this robustness as a premise for
memory-efficient variants:
\begin{itemize}
\item automatic choice of the coreset size~\citep{kamiya2026aso};
\item early-stopped and multi-class lite banks~\citep{lardon2026lite};
\item product-quantised memories for edge deployment~\citep{stropeni2026edge};
\item hierarchical prototype memories with residual normal patches that
restore coverage of under-represented patterns~\citep{feng2026hrpro};
\item density-based coresets motivated by sub-millimetre
defects~\citep{deng2026cpdpatchcore}.
\end{itemize}
Training-free successors refine how the bank is read rather than how large
it is. ProCon~\citep{chae2026procon} replaces nearest-neighbour lookup with
soft projections onto seed-perturbed banks and reports that pixel-level
performance saturates at a 5\,\% budget. Controlled studies of bank
composition compare random, medoid and farthest-first selectors at matched
bank size, but measure contamination rather than defect
size~\citep{chae2026memorycomposition}. Across these works the effect of
compression is measured with aggregate metrics only.

\paragraph{Defect size and evaluation}
The per-region overlap (PRO) curve weights every connected defect region
equally, and its normalised area up to a false-positive rate of
0.3~\citep{bergmann2020uninformed,bergmann2021mvtec} is the standard
localisation metric on MVTec AD. Pixel-level ROC, in contrast, is dominated
by large defects, which motivated the stricter 0.05 limit adopted for MVTec
AD~2~\citep{hecklerkram2026mvtecad2}. Size-conditioned evaluation already
exists:
\begin{itemize}
\item cumulative size-quartile AU-PRO and a size-robustness index, which
show larger size sensitivity on VisA than on MVTec AD for a PatchCore
evaluated at its default memory~\citep{costanzino2025pinpoint};
\item per-image overlap metrics and analyses of region-size
distributions~\citep{bertoldo2024aupimo};
\item synthetic studies of the minimum detectable defect
size~\citep{lehr2024limitations}.
\end{itemize}
Few-shot evaluation on MVTec AD~2, a benchmark with many small defects,
reports that greedy coreset references degrade pixel-level
results~\citep{tu2026fewshotmvtecad2}.
\todo{Read the full texts of \citet{tu2026fewshotmvtecad2} and
\citet{feng2026hrpro} before submission; refine both sentences.}

\paragraph{Protocol variation}
Published numbers on MVTec AD are often not comparable. Implementations
differ in four ways:
\begin{itemize}
\item centre cropping and the resolution at which ground truth is evaluated;
\item how masks are resampled;
\item the PRO normalisation (by the integration limit or by the last
attained false-positive rate);
\item whether PRO is computed over all test images or anomalous images only.
\end{itemize}
EfficientAD~\citep{batzner2024efficientad} re-evaluated competitors without
cropping using the official script. Benchmark critiques document further
inconsistencies~\citep{baitieva2025beyond,rochner2025rethink}. We therefore
fix and report the protocol explicitly (Section~\ref{sec:protocol}).

\paragraph{Gap}
No study we found conditions the effect of memory compression on defect
size at matched memory bytes, with selector controls and seed replication.
This is the gap the present study addresses; its search log is in the
repository (\texttt{docs/LITERATURE.md}).

\section{Experimental protocol}\label{sec:protocol}
\todo{Phase 1: P-REPRO and P-STUDY definitions (DECISIONS D-004/D-005),
pre-registered strata (REALITY\_CHECK H-01), byte accounting (D-012).}

\begin{table}[t]
\centering
\caption{MVTec AD as profiled by our pipeline (generated from manifests).
Median defect area is the per-image native-resolution fraction.}
\label{tab:dataset}
\small
% AUTO-GENERATED by scripts/make_tables.py. DO NOT EDIT.
\begin{tabular}{lrrrrr}
\toprule
Category & Train & Test good & Test defect & Types & Median area (\%) \\
\midrule
\textcolor{red}{[pending]} & & & & & \\
\bottomrule
\end{tabular}

\end{table}

\section{Method under study}\label{sec:method}
\todo{Phase 1: PatchCore as reproduced; selectors; memory accounting.}

\section{Results}\label{sec:results}
\todo{Phase 2: one subsection per supported claim. Evaluator sanity:
constant AU-PRO \SanityConstantAUPRO, oracle AU-PRO \SanityOracleAUPRO.}

\section{Limitations}\label{sec:limitations}
\todo{Single benchmark; MVTec AD reported as comparatively size-stable;
training-free memory detectors only.}

\section{Conclusion}\label{sec:conclusion}
\todo{Written last.}

\section*{Data and code availability}
MVTec AD is available from MVTec Software GmbH under CC BY-NC-SA 4.0 and is
not redistributed. Code, configurations, download and verification scripts,
and the scripts that generate every figure and table are available at
\url{https://github.com/meteorboyF/Anomaly-memory-compression}
\todo{match the journal's required wording}.

\bibliographystyle{elsarticle-harv}
\bibliography{refs}

\end{document}
